\section{Building and verifying the formalization}\label{sec:verify}

The Lean project is the directory \path{lean/} next to this paper. It uses Lean and Mathlib, both at version v4.34.1
(\path{lean-toolchain}, \path{lakefile.toml}, \path{lake-manifest.json}), and requires the library of the eleventh paper of
the series through a relative path; through it, it uses the libraries of the fourth and fifth papers, the shared
library \texttt{lean-machine} and the other libraries those depend on. The sibling directories of the repository must
therefore stay in place.

\paragraph{Build.} With \texttt{elan} installed, run in \path{lean/}:
\begin{center}
\texttt{lake exe cache get \&\& lake build}
\end{center}
The first command downloads the compiled Mathlib; the second compiles the library and the libraries it requires. No
heartbeat limit is raised; the only option set is \texttt{autoImplicit false}.

\paragraph{Axioms.} The file \path{PaperAxioms.lean}, generated from the \LaTeX{} sources of this paper by
\path{tools/lean_index.py}, runs \texttt{\#print axioms} on every declaration named in the paper (\cref{sec:leanindex}):
\begin{center}
\texttt{lake env lean PaperAxioms.lean}
\end{center}
Every declaration depends at most on \texttt{propext}, \texttt{Classical.choice} and \texttt{Quot.sound}. No Lean file of
the library uses \texttt{sorry}, \texttt{admit} or \texttt{native\_decide}, or declares an \texttt{axiom}.

\paragraph{The statement.} The target statements are those of the eleventh paper, \lean{MergeDepth.Q91At} and
\lean{MergeDepth.Q91BinaryAt}. They are built from three earlier definitions:
\begin{itemize}
\item the paired incidence graph \lean{DirectedRecovery.pairedIncidenceGraph} of the fifth paper;
\item the paired-distance condition \lean{DirectedRecovery.HasPairedDistanceSeparation} of the fifth paper;
\item storage codes, \lean{MotionDepthCodes.IsStorageCode}, of the fourth paper.
\end{itemize}
The main formal results, \lean{AboveHalf.not_Q91At_of_tower} and \lean{AboveHalf.not_Q91BinaryAt_of_tower}, prove their
negations. The first needs a family of separated checks (\lean{AboveHalf.CheckSystem}) and a tower of quotients carrying
codes of rate $\rho>1/2$ annihilated by the descended checks; the second needs in addition an injective encoding of the
field in $b$ bits with $\rho\log_2\abs\F/b>1/2$. They are stated for an abstract set of coordinates with a group action and a distance, rather than for
the Bruhat--Tits tree, which is not available in Mathlib; \cref{sec:peeling,sec:descent} explain how the tree provides
these data.

\paragraph{Where things are.} \path{Fluid.lean} (\cref{lem:fluid}), \path{LiftSelection.lean} (\cref{def:selection},
\cref{lem:options,lem:rounding} and the deterministic part of \cref{thm:D}), \path{Peeling.lean} (\cref{lem:peeling}),
\path{Tanner.lean} (\cref{lem:valid,lem:mds}), \path{Descent.lean} (\cref{prop:descent} and the symmetrized graphs of
\cref{sec:undirected}),
\path{Incidence.lean} (degrees and distances in paired incidence graphs), \path{BlowUp.lean} (\cref{lem:blowup}),
\path{Main.lean} (the formal parts of \cref{thm:T}, \cref{thm:A} and \cref{cor:Aprime}), \path{Undirected.lean}
(\cref{lem:graphblowup} and the formal part of \cref{cor:U}) and \path{Numbers.lean} (the numerical facts of
\cref{cor:base,thm:A,cor:Aprime}). \path{Main.lean} and \path{Undirected.lean} do not import \path{Peeling.lean},
\path{Tanner.lean}, \path{LiftSelection.lean}, \path{Fluid.lean} or \path{Numbers.lean}: the main theorems take the
family of checks, the annihilation of the codes and the rate as hypotheses, and the composition of the formal parts is
on paper.
